%----------------------------------------------------------------------------------------
% NOTATIE EN FIGUREN VAN DE CURSUS LOGICA EN BEWIJZEN
% Staat in de preambule en niet in een hoofdstuk: ook de lijsten met oplossingen gebruiken dit.
%----------------------------------------------------------------------------------------
% Bewaarde grafiek, geplaatst door maak-pdf.py: bestand en pagina in figuren-cache/
\newcommand{\grafiekbestand}[2]{\includegraphics[page=#2,trim=3bp 3bp 3bp 3bp]{figuren-cache/#1}}
\usetikzlibrary{positioning}


% Zoals stappen, maar met een bredere kolom voor de berekening: voor lange formules
\newenvironment{stappenbreed}{%
  \par\renewcommand{\arraystretch}{1.25}%
  \noindent\begin{tabular}{@{}>{$\displaystyle}p{9.8cm}<{$}@{\hspace{4mm}}>{\RaggedRight\small\itshape\arraybackslash}p{4.9cm}@{}}%
}{\end{tabular}\par}

% Zoals stappen, met eigen kolombreedtes: \begin{stappenvrij}{6.2cm}{3.8cm}. Voor naast een figuur.
\newenvironment{stappenvrij}[2]{%
  \par\renewcommand{\arraystretch}{1.25}%
  \noindent\begin{tabular}{@{}>{$\displaystyle}p{#1}<{$}@{\hspace{3mm}}>{\RaggedRight\small\itshape\arraybackslash}p{#2}@{}}%
}{\end{tabular}\par}

% Tekst of berekening links, figuur rechts: \naast[breedte rechts]{links}{rechts}
\newsavebox{\naastbox}
\newcommand{\naast}[3][]{%
  \par\noindent\sbox{\naastbox}{#3}%
  \begin{minipage}[t]{\dimexpr\linewidth-\wd\naastbox-3mm\relax}\vspace{0pt}#2\end{minipage}\hfill
  \begin{minipage}[t]{\wd\naastbox}\vspace{0pt}\usebox{\naastbox}\end{minipage}\par\addvspace{1mm}}
% Een rij figuren, gecentreerd
\newcommand{\figuren}[1]{\par\addvspace{1mm}{\centering #1\par}\addvspace{1.5mm}}
% Opgave met een grafiek rechts: \begin{metgrafiek}[breedte]{tikz} opgave en deelvragen \end{metgrafiek}
\newsavebox{\grafiekbox}
\newenvironment{metgrafiek}[2][]{%
  \sbox{\grafiekbox}{#2}%
  \begin{minipage}[t]{\dimexpr\linewidth-\wd\grafiekbox-3mm\relax}%
}{%
  \end{minipage}\hfill
  \begin{minipage}[t]{\wd\grafiekbox}\vspace{-2mm}\usebox{\grafiekbox}\end{minipage}\par}

% Getallenverzamelingen
\newcommand{\N}{\mathbb{N}}
\newcommand{\Z}{\mathbb{Z}}
\newcommand{\Q}{\mathbb{Q}}
\newcommand{\R}{\mathbb{R}}
% Logische symbolen: \en, \of, \niet, \dan (implicatie), \asa (equivalentie)
\newcommand{\en}{\wedge}
\newcommand{\of}{\vee}
\newcommand{\niet}{\neg}
\newcommand{\dan}{\Rightarrow}
\newcommand{\asa}{\Leftrightarrow}
% Deelbaarheid: a \deelt b
\newcommand{\deelt}{\mid}
% Venndiagram in de kleur van het hoofdstuk
\tikzset{verz/.style={draw=accent,thick,fill=accent,fill opacity=.10,text opacity=1}}
